\documentclass[11pt,a4paper]{article}
% ── Pacchetti minimi ────────────────────────────────────────────────
\usepackage[T1]{fontenc}
\usepackage[utf8]{inputenc}
\usepackage[italian]{babel}
\usepackage{amsmath,amssymb}
\usepackage{booktabs}
\usepackage{graphicx}
\usepackage[most]{tcolorbox}

% ── Box "Teorema" (stile tesi_finale) ───────────────────────────────
\newtcolorbox{thmbox}{%
  colback=blue!6, colframe=purple!60!black,
  fonttitle=\bfseries, title=Teorema, breakable
}

\pagestyle{empty}
\hfuzz=20pt
\overfullrule=0pt

\begin{document}

\begin{thmbox}
\textbf{Teorema 3.1 (Convergenza in aspettativa).}
Posto $\alpha := 1-\frac{\lambda}{L}$ e $\beta := \frac{\omega}{2L}$, per ogni $k\ge 0$ vale
\[
\mathbb{E}[J(w_k)-J(w_*)] \;\le\; \alpha^k\varepsilon_0
+ \beta a\,\frac{\alpha^k-a^{-k}}{\alpha a-1},
\]
dove il rapporto va inteso nel senso dei limiti se $\alpha a = 1$.
Di conseguenza, posto $\rho := \max\{\alpha,\,1/a\}$, se $a>1$ si ha $\rho<1$ ed esiste una costante $C=C(\varepsilon_0,\omega,\lambda,L,a)$ tale che
\[
\mathbb{E}[J(w_k)-J(w_*)] \le C\rho^k \qquad \forall\, k\ge 0.
\]
\end{thmbox}

\vspace{1.2em}

\begin{center}
  {\large\bfseries Teorema 3.1 vs.\ Nocedal et al.\ (2012)}
\end{center}

\vspace{0.4em}

\noindent{\footnotesize Simboli: $J$ obiettivo, $J_*$ minimo, $\nabla J$ gradiente; $\lambda$ convessità forte, $L$ lisciatezza; $\theta$ tolleranza; $\omega$ bound di varianza; $n_k=|\mathcal{S}_k|$ batch, $n_k=\lceil a^k\rceil$; $\varepsilon$ accuratezza richiesta.}

\vspace{0.3em}

\renewcommand{\arraystretch}{1.55}
\scriptsize
\begin{tabular}{@{}p{2.6cm}p{4.5cm}p{4.5cm}@{}}
\toprule
\textbf{Aspetto} & \textbf{Nocedal et al.\ (2012)} & \textbf{Questa tesi (Teor.\ 3.1)} \\
\midrule
Stima del gradiente
  & $\|\nabla J\|^2 \ge \lambda\,(J-J_*)$ \newline (perde un fattore 2)
  & $\|\nabla J\|^2 \ge 2\lambda\,(J-J_*)$ \newline (convessità forte) \\
Convergenza deterministica
  & $1-(1-\theta)\lambda/(2L)$
  & $1-(1-\theta)\lambda/L$ \\
Ricorrenza in aspettativa
  & $\varepsilon_{k+1}\le(1-\lambda/2L)\varepsilon_k$ \newline $+\ \omega/(2L n_k)$
  & $\varepsilon_{k+1}\le(1-\lambda/L)\varepsilon_k$ \newline $+\ \omega/(2L n_k)$ \\
Metodo
  & induzione (maggiorazioni)
  & soluzione esatta (forma chiusa) \\
Tasso di convergenza $\rho$
  & $1-\lambda/(4L)$
  & $1-\lambda/L$ \\
Intervallo per $a$
  & $1<a\le(1-\lambda/4L)^{-1}$
  & $1<a\le(1-\lambda/L)^{-1}$ \newline (più ampio) \\
Costante $C$
  & $C=M_{\mathrm N}$
  & $C=C(\varepsilon_0,\omega,\lambda,L,a)$, esplicita \\
Complessità per $\varepsilon$
  & gradienti $O(L/\lambda\varepsilon)$; \newline lavoro $O(Lm/(\lambda\varepsilon)\,M_{\mathrm N})$
  & stesso ordine; \newline lavoro $O(Lm/(\lambda\varepsilon)\,M_{\mathrm T})$, $M_{\mathrm T}<M_{\mathrm N}$ \\
\bottomrule
\end{tabular}

\vspace{0.6em}
\noindent{\footnotesize Note: $\Delta_0=J(w_0)-J(w_*)$; $M_{\mathrm N}=\max\{\Delta_0,\,2\omega/\lambda\}$; $M_{\mathrm T}=\max\{\Delta_0,\,\omega/\lambda\}$. Esito: tasso $\lambda/L$ vs.\ $\lambda/(4L)$; complessità invariata.}

\end{document}
